\section{Применение многосторонних конфиденциальных вычислений в современных системах}
\label{sec:applications}

\subsection{Приватная аналитика данных и аукционы}

В эстонском проекте Sharemind Министерство образования и науки и Налогово-таможенный департамент распределили сведения об обучении и налоговых выплатах между тремя серверами и совместно исследовали связь занятости студентов с завершением обучения. Согласно описанию Cybernetica от 7 июля 2021 г., анализ выполнялся без раскрытия индивидуальных записей и не выявил связи между работой во время учебы и несвоевременным окончанием вуза, но показал зависимость между уровнем образования и доходом \cite{18}.

\subsection{Финансовые приложения и управление ключами}

При распределенном управлении ключом криптовалютного кошелька закрытый ключ разделяется между несколькими участниками, а подпись транзакции формируется совместно без восстановления секрета в одном месте. Такая пороговая схема устраняет единую точку компрометации и позволяет заранее определить число участников, необходимое для подтверждения операции \cite{06}. Тот же подход применим к другим операциям с секретным материалом, например к совместной подписи документов и хранению резервных копий ключей: компрометация одного участника не раскрывает ключ целиком.

\subsection{Машинное обучение и анализ больших данных}
\label{subsec:ml-applications}

МО обычно требует объединения наборов данных, а МКВ позволяет обучать модели и выполнять инференс без передачи данных в открытом виде \cite{04}. Evans и др. выделяют два сценария \cite{06}. При облачном обслуживании запросов протокол объединяет скрытый запрос пользователя и параметры модели сервера, а предсказание получает только пользователь. При совместном обучении несколько организаций строят модель на объединенных данных, не получая доступа к наборам друг друга; для этого используются секретные доли, а часто и гибридные схемы с гомоморфным шифрованием.

В обзоре \cite{07} интеграция МКВ в МО связывается с появлением доступных библиотек. CrypTen (Facebook Research) \cite{22} автоматически преобразует модели PyTorch в версию с МКВ, поддерживает вычисления на GPU и получестную модель нарушителя; интерфейс повторяет API стандартных библиотек МО, что упрощает внедрение. Cerebro содержит язык описания алгоритмов обучения, оптимизатор и средства аудита: компилятор выбирает наиболее эффективные протоколы МКВ, а политики вычислений (compute policies) и криптографический аудит ограничивают выдачу результатов и позволяют проверять корректность вычислений без раскрытия данных \cite{07}.

В работе Mohassel и Zhang \cite{23} данные владельцев разделяются между двумя серверами, которые не вступают в сговор. Протокол обучает линейную и логистическую регрессию методом стохастического градиентного спуска: на каждой итерации веса \(w_j\) обновляются по правилу

\begin{equation}
w_j := w_j - \alpha \cdot \frac{\partial C_i(w)}{\partial w_j},
\end{equation}

где \(\alpha\) --- шаг обучения, \(C_i\) --- функция потерь для выбранного образца, \(\partial C_i(w) / \partial w_j\) --- частная производная по параметру \(w_j\). Для логистической регрессии к линейному слою добавляется сигмоидальная активация; нейронные сети строятся как обобщение таких слоев. Вещественные числа представляются в формате с фиксированной запятой, сложение и умножение выполняются над аддитивными долями, а данные для умножения (тройки Бивера) готовятся заранее на подготовительном этапе. Сигмоида заменяется кусочно-линейной функцией, которая эффективнее вычисляется в МКВ:

\begin{equation}
f(x) =
\begin{cases}
0, & x < -\tfrac{1}{2},\\
x + \tfrac{1}{2}, & -\tfrac{1}{2} \le x \le \tfrac{1}{2},\\
1, & x > \tfrac{1}{2}.
\end{cases}
\end{equation}

Схема не требует доверенного посредника, поскольку данные остаются разделенными между серверами; безопасность доказывается в получестной модели нарушителя. Ограничение схемы --- допущение об отсутствии сговора двух серверов.

Протокол ABY3 \cite{20} рассчитан на три сервера и предполагает, что нарушен не более чем один из них (честное большинство). Значения хранятся в реплицированном аддитивном разделении, поэтому каждый сервер видит лишь часть долей. Протокол допускает преобразования между арифметическим, бинарным представлением и искаженными схемами, содержит усечение для чисел с фиксированной запятой и поддерживает обучение линейной и логистической регрессии и нейронных сетей. Для него описаны варианты, устойчивые к получестному и к злонамеренному нарушителю. TFEncrypted \cite{21} --- библиотека с интерфейсом TensorFlow, в которой ABY3 доступен как один из протоколов для обучения и инференса. Разработчики позиционируют ее как экспериментальное программное обеспечение, поэтому ее использование требует отдельной проверки реализации.

Zhang и др. \cite{12} предлагают протокол для глубоких нейронных сетей на основе порогового деления Шамира. Вещественные числа представляются в формате с фиксированной запятой, а специальный алгоритм усечения ограничивает рост числа дробных разрядов при умножении; это позволяет вычислять свертки и функцию softmax. Умножения оптимизируются алгоритмом Винограда, что, по оценке авторов, сокращает число мультипликативных вентилей более чем на 50\,\%. По сравнению с аддитивным разделением схема эффективнее на подготовительном этапе и при умножении в онлайн-фазе, но требует процедур усечения и поддержки формата с фиксированной запятой.

Издержки на коммуникацию, аппроксимацию нелинейных функций и глубину схем из-за роста шума рассмотрены в разделе~\ref{sec:limitations}. Развитие библиотек направлено на автоматический выбор протоколов и ускорение вычислений на GPU (см. описание платформ выше).

\subsection{Медицина и биоинформатика}

При совместном обучении диагностической модели клиниками данные пациентов остаются локальными, а МКВ защищает вычисление агрегированных обновлений. Обзор \cite{11} относит этот сценарий к основным направлениям конфиденциального медицинского МО; применимость определяется приватностью, вычислительными издержками и масштабируемостью. Медицинские данные относятся к особо чувствительным, а их объединение между организациями ограничено требованиями законодательства о персональных данных и врачебной тайне, поэтому защита на уровне вычислений не заменяет правового основания обработки (см. подраздел~\ref{subsec:standardization}). Дополнительные затраты создают большие объемы данных, например изображений, и разнородность форматов в разных клиниках (см. раздел~\ref{sec:limitations}).

\subsection{Блокчейн и децентрализованные приложения}

В протоколе Bao и др. для блокчейн- и облачных систем многоключевое гомоморфное шифрование позволяет выполнять вычисления над шифротекстами, а направленная расшифровка открывает результат только авторизованной стороне. Подход защищает данные и промежуточные значения без требования взаимного доверия участников \cite{05}. Необходимость такой защиты связана с тем, что данные в блокчейне доступны всем участникам сети, поэтому смарт-контракты не могут напрямую обрабатывать конфиденциальные значения. МКВ позволяют выполнять вычисления над закрытыми данными вне цепочки блоков, а в сам блокчейн, как правило, записывать только результат. Практическими ограничениями остаются задержка и стоимость вычислений и обмена между участниками (см. раздел~\ref{sec:limitations}).

\subsection{Прочие применения: голосование, доказательства с нулевым разглашением и мониторинг}

В электронном голосовании МКВ позволяет конфиденциально суммировать бюллетени; в церемониях формирования параметров доказательств с нулевым разглашением --- распределять генерацию секрета так, чтобы ни один участник не владел им полностью; в сетевом мониторинге --- совместно выявлять совпадения и аномалии без объединения открытых журналов. Во всех сценариях раскрывается только согласованный результат, а исходные данные остаются закрытыми \cite{06}.

\subsection{Вывод}

Наиболее разработанным применением МКВ в рассмотренных сценариях является МО, в частности обучение регрессионных моделей и нейронных сетей на разделенных данных. Существующие решения различаются числом серверов (два в SecureML, три в ABY3 и TFEncrypted) и допущениями о нарушителе: для получестной модели они наиболее зрелы, а защита от злонамеренного нарушителя требует дополнительных затрат. Нелинейные функции, в первую очередь сигмоида, заменяются кусочно-линейными аппроксимациями, а вещественные числа представляются в формате с фиксированной запятой с процедурой усечения. Эти особенности определяют выбор платформы, протокола и состава серверов для MVP, описанного в разделе~\ref{sec:future-research}. Ограничения внедрения МКВ и способы их преодоления рассмотрены в разделе~\ref{sec:limitations}.